\begingroup\small
L’œuvre originale est d’Alexandre Grothendieck : *Fondements de la géométrie algébrique*, exposés Bourbaki 149, 182, 190, 195, 212, 221, 232 et 236, leurs errata et les Commentaires de 1962. \href{https://www.numdam.org/}{Textes originaux : NUMDAM} et \href{https://webusers.imj-prg.fr/~leila.schneps/grothendieckcircle/FGA.pdf}{Grothendieck Circle}.

La transcription moderne, les corrections, l’édition et la composition sont produites par IA. L’œuvre originale conserve son auteur ; elle n’est pas attribuée à l’IA. Le responsable de la publication n’a pas traduit le texte.

Les droits sur les œuvres originales et sur toute contribution protégée à ces traductions et transcriptions dérivées restent ceux de leurs titulaires respectifs. Tout ce sur quoi l'éditeur détient des droits dans ces éditions est dédié au domaine public sous CC0 1.0 (https://creativecommons.org/publicdomain/zero/1.0/deed.fr). Les mentions de droits de tiers et les autorisations déjà accordées restent applicables.

Ces éditions sont mises à disposition gratuitement, à des fins non commerciales et exclusivement pédagogiques, dans la présomption que les titulaires de droits ne s'opposeraient pas à la libre disponibilité de ces traductions et transcriptions aux côtés des originaux français sur \href{https://www.numdam.org/journals/PMIHES/}{NUMDAM}. Cette présomption ne signifie pas qu'une autorisation a été obtenue ; le libre accès aux originaux n'autorise pas, à lui seul, leur traduction, leur adaptation ou leur republication.

Si vous détenez des droits sur un élément présenté ici et vous opposez à sa disponibilité, veuillez \href{https://github.com/KokunoYumeto/fga-fr/issues}{ouvrir une demande sur GitHub}. Le responsable de la publication retirera l'élément concerné à la demande du titulaire des droits.

\par\endgroup\clearpage
